\begin{abstract}
In a one-period CARA economy with Gaussian payoffs, exogenous quantities and
endogenous prices, each investor compares his wealth to a weighted average over
a network of peers. Equilibrium demand scales solve a Katz--Bonacich system, and
prices are explicit in one scalar, the effective market coefficient, which is
itself independent of prices; the market Sharpe ratio is proportional to it. Two
limits, available for every network, organise how the network aggregates risk
aversion. At zero comparison intensity the marginal effect is a quadratic form
in risk tolerances, and a doubly stochastic influence matrix --- balanced
attention --- makes it a Dirichlet form, hence non-negative for every profile of
tolerances. Under intense comparison the coefficient converges to a stationary
average of individual coefficients, weighted by influence rather than by
population. It then falls below its no-comparison value, taking the Sharpe ratio
with it, whenever stationary influence sufficiently favours the less risk
averse; it interpolates from the harmonic to the arithmetic mean when influence
is spread evenly; and it does not move at all when the network splits into
blocks homogeneous in type. A
separate result is that comparison is exactly neutral, for every network and
every profile of intensities, when net risk aversion is common.
\end{abstract}

\noindent\textbf{Keywords:} relative wealth concerns, keeping up with the
Joneses, CARA aggregation, Katz--Bonacich centrality, homophily.

\medskip
\noindent\textbf{JEL:} G11, G12, D53, D85.

%=====================================================================
\section{Introduction}

\cite{gomez} studies an exchange economy in which agents evaluate consumption
relative to a weighted average of the consumption of others, and asks when
equilibrium prices depend only on aggregates. His Proposition~1 shows that
linear risk tolerance with common cautiousness delivers a representative agent
and rules out portfolio bias; his Proposition~2 gives the conditions under
which prices depend on the aggregate endowment alone, and finds them
restrictive: the preference parameters must be common across agents and, absent
the trivial case of uniform weights, so must the equilibrium value of the
endowment. Under those conditions the CARA case reduces cleanly, and keeping up
with the Joneses is equivalent to lowering the representative agent's risk
aversion from $\alpha$ to $\alpha(1-\gamma)$.

What happens outside those conditions is the natural question, and \cite{gomez}
leaves it open. He notes that in their absence equilibrium prices may be
affected by the distribution of weights in the definition of the Joneses, by the
preference parameters, and by the endowment distribution --- three channels that
a representative-agent reduction cannot express. This paper opens the first of
them, in the CARA case, by letting the benchmark be agent-specific: investor $a$
compares himself to $\sum_jA_{aj}W_j$, a weighted average over a network of
influence rather than a single economy-wide aggregate. The question is then not
whether the comparison matters but \emph{what feature of the network} decides
how it matters.

Two remarks on what is being held fixed, since they govern the reading of
everything below. First, with preferences
$-\exp[-\theta_aW_a-\lambda_a(W_a-\sum_jA_{aj}W_j)]$, the coefficient of
absolute risk aversion in own wealth is $\alpha_a=\theta_a+\lambda_a$, not
$\theta_a$. We call $\theta_a$ \emph{net} risk aversion: it is the coefficient
on a common translation of all wealths, and it is risk aversion net of the
comparison. Raising $\lambda_a$ at fixed $\theta_a$ therefore raises own-wealth
risk aversion too, and the comparative statics below are taken along that path.
Second, the economy is closed: payoffs and quantities are the primitives,
prices are endogenous, and Section~\ref{sec:returns} shows that equilibrium
prices are explicit in the primitives and in a single scalar $\bth$ which is
itself independent of prices. The market Sharpe ratio equals $\bth$ times a
constant on the payoff side, so the object characterised below is a price of
risk in the strict sense.

The demand scales solve $[\D_{\theta+\lambda}-\D_{\lambda}A]\vec c=\mathbf 1$, a
Katz--Bonacich system in the sense of \cite{ballester}, and the paper's
contribution is to characterise how that system aggregates. Two limits, both
available for every network, do the work.

At the origin, the marginal effect of comparison is the quadratic form
$\bth'(0)=m\,\vec t^{\top}(\mathbb{I}-A)\vec t/(\mathbf 1^{\top}\vec t)^{2}$ in
the vector $\vec t$ of risk tolerances (Theorem~\ref{thm:deriv}). If attention
is balanced --- $A$ doubly stochastic, so that each investor receives as much
attention as he gives, and not necessarily symmetric --- this is a Dirichlet
form, $\tfrac12\sum_{a,j}A_{aj}(t_a-t_j)^{2}$, hence non-negative and zero only
if tolerance is constant along every link. Comparison raises the effective
coefficient exactly to the extent that investors look at people unlike
themselves.

At the other end, as comparison becomes intense, $\bth$ converges to a
\emph{stationary} average of individual coefficients: $\vec\pi^{\top}\vec\theta$
when the influence matrix has a single recurrent class, and a harmonic
combination across classes in general (Theorem~\ref{thm:intense}). Influence,
not population, does the weighting in the limit, and since the no-comparison
value is the harmonic mean $H(\theta)$, a sufficient condition for comparison to
\emph{lower} the effective coefficient is $\vec\pi^{\top}\vec\theta<H(\theta)$.

These two limits organise what would otherwise be a list. The mean-field
regime is the case in which the stationary distribution is uniform, so the
limit is $\E[\theta]$ and comparison interpolates from the harmonic to the
arithmetic mean (Proposition~\ref{prop:mf}). Sign reversal is the case in which
stationary influence favours the less risk averse; in the three-investor example
of Proposition~\ref{prop:reversal} the most risk averse investor is watched by
nobody, receives stationary weight zero, and the limit falls to $3/4$ against a
starting value of $6/7$. Perfect homophily is the case of several recurrent
classes each homogeneous in type, where the stationary average reproduces the
harmonic mean and the two limits coincide, so the aggregate never moves at all
(Proposition~\ref{prop:homophily}).

A separate and much simpler result delimits where any of this can happen. If net
risk aversion is common, the comparison is exactly neutral: every investor holds
the position he would hold with no relative concern, for every row-stochastic
$A$ and every profile $\lambda_a$ (Theorem~\ref{thm:neutral}). No statistic of
the network survives into equilibrium. This is the network extension of Gómez's
CARA case, and it says that dispersion of net risk aversion is the necessary
fuel for everything else.

The message is that the comparison channel in this class of models is a property
neither of preferences alone nor of the network alone, but of the alignment
between them --- specifically, of the correlation between net risk aversion and
stationary influence. A representative-agent treatment delivers the mean-field
case, and can therefore be wrong in sign, not merely in magnitude, whenever
attention is unbalanced.

%=====================================================================
\section{Related literature}

Preferences over relative consumption or wealth have been used to address a
range of asset-pricing questions. \cite{abel} and \cite{gali} build
consumption-based models with external benchmarks and obtain an amplified
equity premium; \cite{campbell} show how an external habit generates
time-varying effective risk aversion; \cite{chankogan} make the wealth
distribution a state variable under heterogeneous CRRA preferences, which is
the dynamic counterpart of the configuration studied here. \cite{dupor}
provide the taxonomy under which \eqref{eq:pref} is jealousy combined with
keeping up with the Joneses.

Within the CARA class, \cite{degennaro} is the closest antecedent and the
comparison locates what is new here. They study two agents with multiple risky
assets and distributions in a natural exponential family, and show that relative
concerns act through a transformation of effective risk tolerances, lowering
equilibrium premia. The mechanism they identify is a change in the tolerance
\emph{levels}. With two agents the benchmark of each is the other, so there is
no network to speak of: the influence matrix is forced to
$A=\big(\begin{smallmatrix}0&1\\1&0\end{smallmatrix}\big)$, which is doubly
stochastic, and by Theorem~\ref{thm:deriv} the marginal effect is then
necessarily non-negative. The questions of this paper --- whether attention is
balanced, which investors stationary influence favours, whether types and
positions are correlated --- cannot be posed at $m=2$, and the results that turn
on them (Theorems~\ref{thm:deriv} and \ref{thm:intense}, and the reversal and
homophily that follow) have no two-agent counterpart. What is added is therefore
not a dimension count but the object doing the aggregating: a stationary
distribution in place of a pair of tolerances. \cite{demarzo04} obtain portfolio bias when
the benchmark differs in composition from the investor's own holdings, a
channel excluded here by the scalar structure of \eqref{eq:foc}.
\cite{lacker} solve relative-performance games for CARA and CRRA investors in
mean-field and $n$-agent form; their concern is the existence and structure of
Nash equilibria rather than whether the aggregate moves, and their mean-field
formulation is the analogue of our Proposition~\ref{prop:mf} in a different
setting. \cite{kraft} take the problem to incomplete markets in continuous
time. \cite{panageas} surveys what heterogeneity and inequality imply for asset
prices more broadly.

Two connections deserve care. \cite{gollier} gives the criterion governing the
effect of endowment inequality on the equity premium, namely concavity of risk
tolerance in wealth. Our \eqref{eq:cv2} is an analogue for dispersion in a
preference parameter rather than in endowments, and should be read as such
rather than as an application of his theorem. And \cite{ballester} give the
linear-quadratic network game whose solution \eqref{eq:system} reproduces; what
the present setting adds is that the centrality vector becomes constant, and
every statistic of the network drops out, precisely when net risk aversion is
common.

%=====================================================================
\section{The economy}

The primitives are payoffs and quantities, not return moments. There are two
dates. At the second date $n$ risky assets pay $\vec X\in\mathbb{R}^n$, Gaussian
with mean $\vec\mu$ and positive definite covariance $V$; they are in exogenous
supply of $\vec\nu>0$ shares and trade at the first date at prices
$\vec p\in\mathbb{R}^n$, to be determined. A riskless asset pays gross return
$R_f>0$. There are $m$ investors; investor $a$ begins with wealth $W_a^0$ and
chooses share holdings $\vec z_a\in\mathbb{R}^n$, placing the remainder in the
riskless asset, so that
%
\begin{equation}
W_a=R_f\big(W_a^0-\vec p^{\top}\vec z_a\big)+\vec X^{\top}\vec z_a .
\label{eq:wealth}
\end{equation}
%
Preferences are
%
\begin{equation}
U_a=-\exp\Big[-\theta_a W_a-\lambda_a\Big(W_a-\sum_jA_{aj}W_j\Big)\Big],
\qquad \theta_a>0,\ \lambda_a\geq0 .
\label{eq:pref}
\end{equation}

\begin{assumption}[Normalised benchmark]
\label{ass:norm}
$m\geq2$; $A_{aa}=0$, $A_{aj}\geq0$ and $\sum_jA_{aj}=1$ for every $a$.
\end{assumption}

\begin{definition}[Two coefficients]
\label{def:two}
Holding the wealth of others fixed, the coefficient of absolute risk aversion
of \eqref{eq:pref} in own wealth is
%
\begin{equation}
\alpha_a=\theta_a+\lambda_a ,
\label{eq:alpha}
\end{equation}
%
which we call \emph{own-wealth risk aversion}. We call $\theta_a$ the
\emph{net risk-aversion coefficient}: under Assumption~\ref{ass:norm} it is the
coefficient on a common translation of all wealths, since adding a constant to
every $W_j$ shifts the exponent of \eqref{eq:pref} by $-\theta_a$ times that
constant, independently of $\lambda_a$. It is risk aversion net of the
comparison.
\end{definition}

The distinction matters for reading everything below, and we state it before
any result because the two coefficients cannot both be held fixed. Raising
$\lambda_a$ at constant $\theta_a$ is not a pure increase in the intensity of
comparison: by \eqref{eq:alpha} it raises own-wealth risk aversion by the same
amount. Every comparative static in $\lambda$ in this paper is taken along that
path, at constant $\theta$, and the neutrality of Section~\ref{sec:neutral} is
exactly the cancellation of the two effects. The alternative path, at constant
$\alpha$, is described in Remark~\ref{rem:param}.

The multiplicative parametrization of \cite{gomez} makes the same point in his
notation. Collecting terms in the exponent of \eqref{eq:pref},
%
\begin{equation}
-\theta_aW_a-\lambda_a\Big(W_a-\sum_jA_{aj}W_j\Big)
=-\alpha_a\Big(W_a-\gamma_a\sum_jA_{aj}W_j\Big),
\qquad \gamma_a=\frac{\lambda_a}{\theta_a+\lambda_a}\in[0,1),
\label{eq:param}
\end{equation}
%
so that $\theta_a=\alpha_a(1-\gamma_a)$. He writes the CARA case as
$u(c,X)=-\exp(-\alpha(c-\gamma X))$ and reduces it to a representative agent
with coefficient $\alpha(1-\gamma)$, which is exactly $\theta$.

The benchmark is thus a weighted average of the wealth of the other investors,
as in \cite{gomez}, with the difference that the weights may vary across
investors. Assumption~\ref{ass:norm} is what gives $\theta_a$ the reading in
Definition~\ref{def:two}, and it is what the results below require. Dropping it changes the object
being compared: row sums below one describe comparison with a partial or
external benchmark, which is a coherent alternative we do not study, while row
sums above $1+\theta_a/\lambda_a$ make an investor worse off when the whole
economy becomes richer, which is harder to motivate. We retain
Assumption~\ref{ass:norm} throughout and claim nothing about either case.

\begin{assumption}[Nash behaviour]
\label{ass:nash}
Each investor chooses $\vec z_a$ taking the holdings $\vec z_j$, $j\neq a$,
as fixed at their equilibrium values. An equilibrium is a profile
$(\vec z_1,\dots,\vec z_m)$ and a price vector $\vec p$ such that each
$\vec z_a$ is optimal given the others and markets clear.
\end{assumption}

Assumption~\ref{ass:nash} is the standard notion and we state it explicitly
because the conjecture it embodies matters for what follows. An investor
behaving in this way conjectures that the positions of others do not respond to
his own: the conjectural derivative is zero, and the conjectured levels are the
equilibrium positions.

\subsection{First-order conditions and equilibrium prices}

Substituting \eqref{eq:wealth} into \eqref{eq:pref}, the random part of the
exponent is $-\vec X^{\top}\vec q_a$ with
%
\begin{equation}
\vec q_a=(\theta_a+\lambda_a)\vec z_a-\lambda_a\sum_jA_{aj}\vec z_j ,
\label{eq:qa}
\end{equation}
%
and the deterministic part is
$-(\theta_a+\lambda_a)R_f(W_a^0-\vec p^{\top}\vec z_a)
+\lambda_a\sum_jA_{aj}R_f(W_j^0-\vec p^{\top}\vec z_j)$. Gaussian payoffs give
$\E[U_a]=-\exp\big(d_a-\vec q_a^{\top}\vec\mu+\tfrac12\vec q_a^{\top}V\vec q_a\big)$
with $d_a$ the deterministic part. Since
$\partial\vec q_a/\partial\vec z_a=(\theta_a+\lambda_a)\mathbb{I}$, the exponent
is a strictly convex quadratic in $\vec z_a$, with Hessian
$(\theta_a+\lambda_a)^{2}V$. The optimum is therefore interior and solves
%
\begin{equation}
V\vec q_a=\vec\mu-R_f\vec p .
\label{eq:foc}
\end{equation}
%
Every solution is proportional to $V^{-1}(\vec\mu-R_f\vec p)$, so all investors
hold the same portfolio of risky assets and differ only in scale; two-fund
separation holds, as in Gómez's Proposition~1. Writing
$\vec z_a=c_aV^{-1}(\vec\mu-R_f\vec p)$ reduces \eqref{eq:foc} to
%
\begin{equation}
\big[\D_{\theta+\lambda}-\D_{\lambda}A\big]\,\vec c=\mathbf 1,
\qquad \D_x=\mathrm{diag}(x_1,\dots,x_m),
\label{eq:system}
\end{equation}
%
whose solution is a Katz--Bonacich centrality of the influence network
\cite{ballester}, generalised in that the attenuation factors $\lambda_a$ and
the base vector are agent-specific rather than scalar. Write
$M=\D_{\theta+\lambda}-\D_{\lambda}A$. Its off-diagonal entries are
non-positive, so $M$ is a Z-matrix, and its off-diagonal row sums are
$\lambda_a\sum_jA_{aj}=\lambda_a<\theta_a+\lambda_a$, so it is strictly
diagonally dominant with positive diagonal. A strictly diagonally dominant
Z-matrix with positive diagonal is a non-singular M-matrix, hence $M^{-1}\geq0$
entrywise \cite{bermanplemmons}; existence and uniqueness follow, and
$\vec c=M^{-1}\mathbf 1>0$ because no row of $M^{-1}$ can vanish. Define
%
\begin{equation}
\frac{1}{\theta_M}=\sum_ac_a,
\qquad \bth=m\,\theta_M ,
\label{eq:thm}
\end{equation}
%
and call $\bth$ the effective market coefficient. Market clearing
$\sum_a\vec z_a=\vec\nu$ then gives $\vec\mu-R_f\vec p=\theta_MV\vec\nu$, that
is
%
\begin{equation}
\vec p=\frac{\vec\mu-\theta_M\,V\vec\nu}{R_f} .
\label{eq:price}
\end{equation}
%
Equation \eqref{eq:price} closes the model. Prices are explicit in the
primitives $(\vec\mu,V,\vec\nu,R_f)$ and in the single scalar $\theta_M$, and no
fixed point has to be solved, because $\theta_M$ is determined by
\eqref{eq:system} and \eqref{eq:thm} from $(\vec\theta,\vec\lambda,A)$ alone and
does not depend on prices. This is a property of CARA preferences with Gaussian
payoffs, not an approximation: demands are independent of initial wealth, so
the wealth distribution never feeds back. Prices are positive, and the
equilibrium admissible, exactly when
%
\begin{equation}
\vec\mu>\theta_M\,V\vec\nu
\qquad\text{componentwise},
\label{eq:positive}
\end{equation}
%
which is the only parameter restriction the closure imposes; parameters
satisfying it are called admissible. If in addition $\vec\mu\gg0$, every
sufficiently small $\theta_M$ is admissible. All the results below concern $\theta_M$, so
they are statements about \eqref{eq:price}.

\subsection{Returns and the price of risk}
\label{sec:returns}

Gross returns are $R_i=X_i/p_i$. Writing $\D_p=\mathrm{diag}(\vec p)$, the
return covariance and the vector of expected excess returns are
%
\begin{equation}
\Sigma=\D_p^{-1}V\D_p^{-1},
\qquad
\vec r_e=\D_p^{-1}\big(\vec\mu-R_f\vec p\big) ,
\label{eq:returns}
\end{equation}
%
both endogenous. Let $W_M=\vec p^{\top}\vec\nu$ be the market value of the risky
assets and $\vec w=\D_p\vec\nu/W_M$ its vector of value weights. Substituting
\eqref{eq:price} into \eqref{eq:returns} gives
%
\begin{equation}
\vec r_e=\theta_M\,W_M\,\Sigma\vec w ,
\label{eq:clear}
\end{equation}
%
the mean--covariance relation, now with every object in it determined. The
market portfolio has excess return $R_M-R_f$ with
$\E[R_M]-R_f=\theta_M\vec\nu^{\top}V\vec\nu/W_M$ and
$\mathrm{Var}(R_M)=\vec\nu^{\top}V\vec\nu/W_M^{2}$, so
%
\begin{equation}
\frac{\E[R_M]-R_f}{\mathrm{Var}(R_M)}=\theta_MW_M ,
\qquad
S_M=\frac{\E[R_M]-R_f}{\mathrm{sd}(R_M)}
=\frac{\bth}{m}\sqrt{\vec\nu^{\top}V\vec\nu} .
\label{eq:sharpe}
\end{equation}
%
The first is the market price of risk per unit of variance; the second is the
Sharpe ratio, and it is $\bth$ multiplied by a constant that lives entirely on
the payoff side. Every comparative static in $\bth$ below is therefore a
comparative static in the price of risk, in the strict sense. The
unconditional statements are the aggregate ones: a fall in $\bth$ lowers the
market Sharpe ratio and the market risk premium. Statements about individual
prices are conditional, since
$\partial p_i/\partial\theta_M=-(V\vec\nu)_i/R_f$, so that a fall in $\bth$
raises $p_i$ exactly when $(V\vec\nu)_i>0$ --- when asset $i$ covaries
positively with the market payoff. Assets that hedge the market move the other
way.

\begin{remark}[On the riskless rate]
\label{rem:rf}
Taking $R_f$ as given amounts to a riskless asset in perfectly elastic supply,
and then $W_M=\vec p^{\top}\vec\nu$ is an endogenous market value that need not
equal aggregate initial wealth. If instead the riskless asset is in zero net
supply, the budget identities give $\sum_a(W_a^0-\vec p^{\top}\vec z_a)
=\sum_aW_a^0-W_M$, so clearing that market requires
$W_M=\sum_aW_a^0$ and determines
%
\begin{equation*}
R_f=\frac{\vec\mu^{\top}\vec\nu-\theta_M\,\vec\nu^{\top}V\vec\nu}
{\sum_aW_a^0} .
\end{equation*}
%
Both closures leave \eqref{eq:price}--\eqref{eq:sharpe} intact, since neither
touches $\theta_M$; they differ only in whether $W_M$ or $R_f$ is the free
quantity. Under the second, the market value of the risky assets equals
aggregate initial wealth, and since
$\partial R_f/\partial\theta_M=-\vec\nu^{\top}V\vec\nu/\sum_aW_a^0<0$,
comparison \emph{raises} the riskless rate whenever it lowers $\bth$. That
closure also fixes $\vec\nu^{\top}\vec p=\sum_aW_a^0$, so prices cannot all move
in the same direction: a fall in $\bth$ redistributes value across assets
instead of raising the aggregate.
\end{remark}

%=====================================================================
\section{Neutrality under common net risk aversion}
\label{sec:neutral}

\begin{theorem}
\label{thm:neutral}
Let Assumptions~\ref{ass:norm} and \ref{ass:nash} hold and suppose
$\theta_a=\theta$ for every $a$, that is, $\alpha_a(1-\gamma_a)=\theta$. Then $\vec c=\theta^{-1}\mathbf 1$ is the
unique solution of \eqref{eq:system}, so every investor holds
$\vec z_a=V^{-1}(\vec\mu-R_f\vec p)/\theta$ and $\bth=\theta$, for every influence
matrix satisfying Assumption~\ref{ass:norm} and every profile
$(\lambda_1,\dots,\lambda_m)$.
\end{theorem}

\begin{proof}
Substituting $\vec c=\theta^{-1}\mathbf 1$ in row $a$ of \eqref{eq:system} and
using $\sum_jA_{aj}=1$ gives
$(\theta+\lambda_a)/\theta-\lambda_a/\theta=1$. Uniqueness follows from strict
diagonal dominance, as above.
\end{proof}

The hypothesis is that the product $\alpha_a(1-\gamma_a)$ be common, so
neutrality holds on the whole level set
$\{(\alpha,\gamma):\alpha(1-\gamma)=\theta\}$. Along that set an investor with a
stronger comparison motive is also more risk averse in own wealth, in exactly
the proportion that cancels. That is the content of the theorem, and it is
weaker than it sounds if $\theta_a$ is read as the conventional coefficient of
absolute risk aversion, which by Definition~\ref{def:two} it is not.

The cancellation occurs row by row, before any aggregation, which is why
nothing about $A$ or about the dispersion of $\lambda_a$ enters. Economically, the
benchmark is fixed in a unilateral deviation: under Assumption~\ref{ass:nash}
the positions of others do not respond. Neutrality arises instead because at
the equilibrium profile the benchmark carries the same risky exposure as the
investor himself, so that
$(\theta+\lambda_a)\vec z-\lambda_aA\vec z=(\theta+\lambda_a)\vec z-\lambda_a\vec z=\theta\vec z$.
The benchmark exposure offsets the additional $\lambda_a$ loading on own
equilibrium wealth in the level of total risky exposure, not through any
conjectured response to the deviation.

Theorem~\ref{thm:neutral} is the network extension of the CARA case of
\cite{gomez}, and we claim no more for it than that. His Proposition~2
identifies common cautiousness --- here, common $\theta$ --- as what makes the
weighting drop out; the addition is that the weighting may be agent-specific
and the comparison intensities may differ, and neither disturbs the conclusion.

\begin{remark}[Comparison at constant own-wealth risk aversion]
\label{rem:param}
Theorem~\ref{thm:neutral} varies $\lambda_a$ holding $\theta$ fixed. The other
natural path holds $\alpha$ fixed and varies $\gamma_a$, so that
$\theta_a=\alpha(1-\gamma_a)$ is dispersed whenever $\gamma_a$ is. Neutrality
then fails, but not through a new mechanism: the demand scales solve
\eqref{eq:system} with $\theta_a=\alpha(1-\gamma_a)$ and
$\lambda_a=\alpha\gamma_a$, and dispersion of the keeping-up weight is
dispersion of net risk aversion. The two paths are therefore two coordinates of
one statement. What neutrality requires is that $\alpha_a(1-\gamma_a)$ be
common; what breaks it is dispersion in that product, however generated.

Gómez's Proposition~2 requires more, in a second way. It asks that the
preference parameters be identical across agents and, barring uniform weights,
that the equilibrium value of the endowment be identical as well. Initial
wealth is absent from \eqref{eq:foc} altogether, since CARA demands have no
wealth effects, so no restriction on the endowment distribution is needed here.
\end{remark}

\begin{remark}[Two notions of neutrality]
\label{rem:two}
Theorem~\ref{thm:neutral} asserts that individual positions are unchanged,
which is stronger than saying the effective coefficient is unchanged. The latter
requires only that $\sum_ac_a$ be unchanged, and deviations can offset. The
distinction is not hypothetical: it occurs inside the example of
Proposition~\ref{prop:reversal}, whose $\bth(\lambda)$ returns to its
no-comparison value at $\lambda=\tfrac13$. There
%
\begin{equation*}
\vec c(0)=\big(2,\;1,\;\tfrac12\big),
\qquad
\vec c(\tfrac13)=\big(\tfrac53,\;\tfrac76,\;\tfrac23\big),
\qquad
\bth(0)=\bth(\tfrac13)=\tfrac67 ,
\end{equation*}
%
so the aggregate is exactly neutral while individual positions move by
$-16.7$, $+16.7$ and $+33.3$ per cent. Results below concern $\bth$, so the
weaker notion is the relevant one, and statements about it should not be read
as statements about allocations.
\end{remark}

%=====================================================================
\section{Heterogeneous net risk aversion}
\label{sec:hetero}

We now let $\theta_a$ vary, which is exactly the configuration in which the
aggregation conditions of \cite{gomez} fail. The system \eqref{eq:system} still
has a unique positive solution, but it is no longer proportional to
$\mathbf 1$, and $\bth$ depends on $\lambda$ and on $A$.

Throughout this section $\lambda_a=\lambda$ is common, so that
$\lambda$ indexes a one-parameter family of economies running from no
comparison to complete comparison. Two limits of that family can be computed in
closed form for \emph{every} network, and between them they organise everything
that follows. Write
%
\begin{equation}
\vec t=\D_{\theta}^{-1}\mathbf 1,
\qquad t_a=1/\theta_a ,
\label{eq:tol}
\end{equation}
%
for the vector of risk tolerances, so that $\vec c(0)=\vec t$ and
$\bth(0)=H(\theta)$ is the harmonic mean of net risk aversion.

\subsection{The marginal effect of comparison}
\label{sec:marginal}

\begin{theorem}
\label{thm:deriv}
Let Assumption~\ref{ass:norm} hold and $\lambda_a=\lambda$. Then
%
\begin{equation}
\bth'(0)=\frac{m\,\vec t^{\top}(\mathbb{I}-A)\vec t}
{\big(\mathbf 1^{\top}\vec t\big)^{2}} .
\label{eq:deriv0}
\end{equation}
%
If in addition $A$ is doubly stochastic --- not necessarily symmetric --- then
%
\begin{equation}
\vec t^{\top}(\mathbb{I}-A)\vec t
=\tfrac12\sum_{a,j}A_{aj}\big(t_a-t_j\big)^{2}\;\geq\;0 ,
\label{eq:dirichlet}
\end{equation}
%
so $\bth'(0)\geq0$, with equality if and only if $t_a=t_j$ whenever $A_{aj}>0$.
For $A$ merely row stochastic the sign of \eqref{eq:deriv0} is unrestricted.
\end{theorem}

\begin{proof}
Write $M(\lambda)=\D_{\theta}+\lambda(\mathbb{I}-A)$, so that
$\vec c=M^{-1}\mathbf 1$ and
$\vec c\,'=-M^{-1}(\mathbb{I}-A)\vec c$. At $\lambda=0$ we have
$M(0)=\D_{\theta}$ and $\vec c(0)=\vec t$, hence
$\mathbf 1^{\top}\vec c\,'(0)=-\mathbf 1^{\top}\D_{\theta}^{-1}(\mathbb{I}-A)\vec t
=-\vec t^{\top}(\mathbb{I}-A)\vec t$, using
$\mathbf 1^{\top}\D_{\theta}^{-1}=\vec t^{\top}$. Since
$\bth=m/(\mathbf 1^{\top}\vec c)$,
$\bth'=-m(\mathbf 1^{\top}\vec c\,')/(\mathbf 1^{\top}\vec c)^{2}$, which gives
\eqref{eq:deriv0}. For the second claim, row and column sums of $A$ both equal
one, so
$\sum_{a,j}A_{aj}(t_a^{2}+t_j^{2})/2
=\tfrac12\big(\sum_at_a^{2}+\sum_jt_j^{2}\big)=\vec t^{\top}\vec t$, and
therefore
$\tfrac12\sum_{a,j}A_{aj}(t_a-t_j)^{2}
=\vec t^{\top}\vec t-\vec t^{\top}A\vec t$, which is \eqref{eq:dirichlet}. For
the last claim take
%
\begin{equation*}
A=\begin{pmatrix}0&1&0\\1&0&0\\0&1&0\end{pmatrix},
\qquad
\vec\theta=(1,1,2),
\qquad
\vec t=\big(1,1,\tfrac12\big),
\end{equation*}
%
which satisfies Assumption~\ref{ass:norm} and has column sums $(1,2,0)$. Here
$A\vec t=\mathbf 1$, so $\vec t^{\top}A\vec t=\tfrac52$ while
$\vec t^{\top}\vec t=\tfrac94$, giving
$\vec t^{\top}(\mathbb{I}-A)\vec t=-\tfrac14$ and
$\bth'(0)=3(-\tfrac14)/(\tfrac52)^{2}=-\tfrac{3}{25}<0$.
\end{proof}

In that example the two least risk averse investors absorb all the attention and
the most risk averse receives none, so $\vec\pi=(\tfrac12,\tfrac12,0)$ and
$\vec\pi^{\top}\vec\theta=1<\tfrac65=H(\theta)$. The effect is negative at the
origin and the limit lies below the starting value, and $\bth$ in fact falls
monotonically from $6/5$ to $1$. It is the complement of
Proposition~\ref{prop:reversal} below, where the marginal effect is positive and
the decline sets in only at higher intensity.

Equation \eqref{eq:dirichlet} is a Dirichlet form: the marginal effect of
comparison is a weighted sum of squared differences in risk tolerance across
the links of the network. It is positive exactly when investors look at
somebody unlike themselves, and the relevant notion of ``unlike'' is difference
in tolerance, not in aversion. Double stochasticity is what makes the quadratic
form a Dirichlet form; economically it says that attention is balanced, each
investor receiving as much of it as he gives. It is the balance of attention,
not the profile of tolerances, that the condition constrains: for a given
$\vec t$ the form may well be non-negative without it, but double stochasticity
is what makes it non-negative for \emph{every} $\vec t$.

\subsection{The limit of intense comparison}
\label{sec:intense}

The other end of the family is governed by the stationary distribution of $A$
read as a Markov chain. Let $C_1,\dots,C_r$ be its recurrent classes,
$\vec\pi^{(k)}$ the stationary distribution on $C_k$ normalised to sum to one,
and $h_k(a)$ the probability that a chain started at $a$ is absorbed into
$C_k$.

\begin{lemma}
\label{lem:bound}
$c_a\leq1/\min_b\theta_b$ for every $a$ and every $\lambda\geq0$.
\end{lemma}

\begin{proof}
Let $c^{*}=c_{a^{*}}$ be the largest entry. Row $a^{*}$ of \eqref{eq:system}
reads $(\theta_{a^{*}}+\lambda)c^{*}=1+\lambda\sum_jA_{a^{*}j}c_j\leq1+\lambda c^{*}$,
so $\theta_{a^{*}}c^{*}\leq1$.
\end{proof}

\begin{theorem}
\label{thm:intense}
Let Assumption~\ref{ass:norm} hold and $\lambda_a=\lambda$. Then
%
\begin{equation}
\lim_{\lambda\to\infty}\bth(\lambda)
=\frac{m}{\displaystyle\sum_{k=1}^{r}
\frac{\mathbf 1^{\top}\vec h_k}{\vec\pi^{(k)\top}\vec\theta}} .
\label{eq:intense}
\end{equation}
%
In particular, if $A$ has a single recurrent class with stationary distribution
$\vec\pi$, then $\bth(\lambda)\to\vec\pi^{\top}\vec\theta$.
\end{theorem}

\begin{proof}
The kernel of $\mathbb{I}-A$ is spanned by the absorption vectors
$\vec h_1,\dots,\vec h_r$, which satisfy $\sum_k\vec h_k=\mathbf 1$. Dividing
$M(\lambda)\vec c=\mathbf 1$ by $\lambda$ gives
$(\lambda^{-1}\D_{\theta}+\mathbb{I}-A)\vec c=\lambda^{-1}\mathbf 1$, and
$\vec c$ is bounded by Lemma~\ref{lem:bound}, so any limit point $\vec c_\infty$
satisfies $(\mathbb{I}-A)\vec c_\infty=0$, that is
$\vec c_\infty=\sum_k\kappa_k\vec h_k$. To identify $\kappa_k$, extend
$\vec\pi^{(k)}$ by zero off $C_k$ and left-multiply $M(\lambda)\vec c=\mathbf 1$
by $\vec\pi^{(k)\top}$. Since $C_k$ is closed, $\vec\pi^{(k)\top}A=\vec\pi^{(k)\top}$,
so the term in $\lambda$ vanishes identically and
$\vec\pi^{(k)\top}\D_{\theta}\vec c=1$ for every $\lambda$. Passing to the limit
and using $h_l=\delta_{kl}$ on $C_k$ gives
$\kappa_k\,\vec\pi^{(k)\top}\vec\theta=1$. Then
$\mathbf 1^{\top}\vec c_\infty=\sum_k\kappa_k\mathbf 1^{\top}\vec h_k$, and
$\bth=m/(\mathbf 1^{\top}\vec c)$ gives \eqref{eq:intense}. The limit point is
unique, so the whole family converges.
\end{proof}

Theorem~\ref{thm:intense} is the organising result of the paper. Under intense
comparison the aggregate is no longer the harmonic mean of net risk aversion
but a \emph{stationary} average: influence, not population, does the weighting.
Reading $\bth(0)=H(\theta)$ against \eqref{eq:intense} gives at once a
sufficient condition for the comparison to lower the effective coefficient,
%
\begin{equation}
\vec\pi^{\top}\vec\theta<H(\theta)
\quad\Longrightarrow\quad
\bth(\lambda)<\bth(0)\ \text{ for all large }\lambda ,
\label{eq:revcond}
\end{equation}
%
in the single-class case. Reversal is therefore not an artefact of small
examples but a statement about the correlation between net risk aversion and
stationary influence: it occurs when the investors who attract attention in the
long run are less risk averse than the population as a whole. The three results
that follow are the three ways \eqref{eq:intense} can turn out.

\subsection{The mean-field limit}

\begin{proposition}
\label{prop:mf}
Let $\lambda_a=\lambda$ for all $a$ and consider a sequence of economies
indexed by $m$ in which the empirical distribution of $\theta_a$ converges
weakly to a law. Suppose the networks satisfy
%
\begin{equation}
\theta_a\geq\underline\theta>0\ \text{ for all }a,\ m,
\qquad
\delta_m\;=\;\max_a\Big\|A_{a\cdot}-\tfrac1m\mathbf 1\Big\|_1\longrightarrow0 .
\label{eq:primitive}
\end{equation}
%
Then, writing $\bth_m$ for the effective
coefficient of the $m$-investor economy,
%
\begin{equation}
\bth_m(\lambda)\longrightarrow \bth_\infty(\lambda)=H(\theta+\lambda)-\lambda,
\qquad H(x)=\big(\E[1/x]\big)^{-1} .
\label{eq:mf}
\end{equation}
%
The limit $\bth_\infty$ is differentiable wherever the moments below exist,
and
%
\begin{equation}
\frac{d\bth_\infty}{d\lambda}
=\frac{\mathrm{Var}\big[(\theta+\lambda)^{-1}\big]}
{\E\big[(\theta+\lambda)^{-1}\big]^{2}}\;\geq\;0 ,
\label{eq:cv2}
\end{equation}
%
with equality if and only if $\theta$ is degenerate. If moreover
$\E[\theta^3]<\infty$ then
$\bth_\infty=\E[\theta]-\mathrm{Var}(\theta)/\lambda+O(\lambda^{-2})$ as
$\lambda\to\infty$. Equation \eqref{eq:cv2} is the derivative of the limit;
we do not claim it is the limit of the derivatives of the finite economies.
\end{proposition}

\begin{proof}
Since $\theta_a\geq\underline\theta$, the variable $(\theta+\lambda)^{-1}$ is
bounded by $(\underline\theta+\lambda)^{-1}$, so weak convergence alone gives
convergence of $\E[(\theta+\lambda)^{-1}]$ and no separate integrability
hypothesis is needed. By Lemma~\ref{lem:bound},
$\|\vec c\|_\infty\leq1/\underline\theta$ uniformly in $m$ and $\lambda$. Hence
$\big|\sum_jA_{aj}c_j-\bar c_m\big|
=\big|(A_{a\cdot}-m^{-1}\mathbf 1)^{\top}\vec c\big|
\leq\delta_m\|\vec c\|_\infty\leq\delta_m/\underline\theta\to0$
uniformly in $a$, which is the mixing condition in endogenous form and is where
the uniform bound is needed. Under it, row $a$ of \eqref{eq:system} becomes
$(\theta_a+\lambda)c_a=1+\lambda\bar c$ in the limit, so
$c_a=(1+\lambda\bar c)/(\theta_a+\lambda)$. Averaging and writing
$M=\E[(\theta+\lambda)^{-1}]$ gives $\bar c=M/(1-\lambda M)$ and
$\bth=1/M-\lambda$. Differentiating gives \eqref{eq:cv2}, whose sign is
Cauchy--Schwarz, with equality only for degenerate $\theta$. The expansion
follows from
$\E[(\theta+\lambda)^{-1}]
=\lambda^{-1}\big(1-\E[\theta]/\lambda+\E[\theta^{2}]/\lambda^{2}
+O(\lambda^{-3})\big)$, which requires the third moment to control the
remainder.
\end{proof}

In this limit comparison interpolates the aggregation of risk aversion between
the two classical aggregators: the harmonic mean at $\lambda=0$, which is the
standard CARA result, and the arithmetic mean as $\lambda\to\infty$. Since
$H(\theta)\leq\E[\theta]$, the largest increase in the effective coefficient
available is the harmonic--arithmetic gap of the net risk-aversion
distribution, a
statistic of the cross-section rather than a free parameter.

Condition \eqref{eq:primitive} is primitive: it constrains only $A$ and the
support of $\theta$. The leave-one-out network $A_{aj}=1/(m-1)$, $j\neq a$, has
$\delta_m=2/m$ exactly, so it qualifies; exact equality at finite $m$ would
require identical rows, which Assumption~\ref{ass:norm} forbids since
$A_{aa}=0$. More generally $\delta_m$ is twice the total-variation distance of
each row from the uniform distribution, so \eqref{eq:primitive} asks that every
investor's benchmark be close to the population average --- which is what
``mean field'' should mean, and is checkable from the network alone. The uniform
lower bound $\underline\theta$ cannot be dropped: it is what converts closeness
of rows into closeness of benchmarks, since $\vec c$ is controlled only through
Lemma~\ref{lem:bound}.

\begin{remark}[Consistency with Theorem~\ref{thm:intense}]
\label{rem:consistency}
Condition \eqref{eq:primitive} forces the stationary distribution towards the
uniform one, and \eqref{eq:intense} then returns
$\vec\pi^{\top}\vec\theta\to\E[\theta]$, which is the $\lambda\to\infty$ end of
\eqref{eq:mf}. The mean-field regime is thus the case of
Theorem~\ref{thm:intense} in which influence is spread evenly, and
$\bth'(0)\geq0$ in it because near-uniform rows are near doubly stochastic, so
\eqref{eq:dirichlet} applies approximately. Both directions of the interpolation
are consequences of the two limits.
\end{remark}

\subsection{Non-monotonicity and sign reversal}
\label{sec:reversal}

\begin{proposition}
\label{prop:reversal}
There exist normalised benchmarks and net risk-aversion profiles for which
$d\bth/d\lambda<0$ on a non-empty interval and for which
$\bth(\lambda)<\bth(0)$ at sufficiently large comparison intensities.
\end{proposition}

\begin{proof}
Take $m=3$, $\theta=(\tfrac12,1,2)$, $\lambda_a=\lambda$ and
%
\begin{equation*}
A=\begin{pmatrix}0&1&0\\1&0&0\\1&0&0\end{pmatrix},
\end{equation*}
%
which satisfies Assumption~\ref{ass:norm}. Solving \eqref{eq:system} gives
%
\begin{equation}
\bth(\lambda)=\frac{3(\lambda+2)(3\lambda+1)}{12\lambda^{2}+24\lambda+7},
\qquad
\bth'(\lambda)=-\frac{3\big(12\lambda^{2}+6\lambda-1\big)}
{\big(12\lambda^{2}+24\lambda+7\big)^{2}} .
\label{eq:cex}
\end{equation}
%
Hence $\bth'(0)=3/49>0$ and $\bth'(\lambda)<0$ for
$\lambda>(\sqrt{21}-3)/12\approx0.132$. The function rises from
$\bth(0)=6/7\approx0.857$ to a maximum $\approx0.860$ at that point, then
falls, returning to $6/7$ at $\lambda=1/3$ and continuing downwards:
$\bth(1)=0.837209$ and $\bth(2)=0.815534$.
\end{proof}

Both features of the example are read off the two limits of
Section~\ref{sec:marginal} and Section~\ref{sec:intense}, and neither depends on
the closed form. The column sums of $A$ are $(2,1,0)$, so $A$ is not doubly
stochastic and \eqref{eq:dirichlet} does not apply; \eqref{eq:deriv0}
nonetheless returns $\bth'(0)=3/49>0$, so comparison initially raises the
effective coefficient. At the other end, investor~3 is transient and
$\{1,2\}$ is the unique recurrent class, with $\vec\pi=(\tfrac12,\tfrac12,0)$.
Theorem~\ref{thm:intense} therefore gives
%
\begin{equation*}
\lim_{\lambda\to\infty}\bth(\lambda)
=\vec\pi^{\top}\vec\theta=\tfrac12\big(\tfrac12+1\big)=\tfrac34
\;<\;\tfrac67=H(\theta)=\bth(0),
\end{equation*}
%
so \eqref{eq:revcond} holds and reversal is forced. The economics is now
explicit. The most risk averse investor watches the least risk averse and is
watched by nobody, so stationary influence assigns him weight zero: his risk
aversion is present in the harmonic mean at $\lambda=0$ and absent from the
stationary mean at $\lambda=\infty$. Non-monotonicity is the collision of the
two limits --- a positive marginal effect at the origin with a strictly lower
value at infinity --- and the closed form only locates the turning point.

How often reversal occurs is a question about a distribution over networks, and
the answer is a property of that distribution rather than of the mechanism.
Under a dense i.i.d.-weight design, specified in the supplement, it becomes rare
as the number of investors grows. That ensemble does not satisfy
\eqref{eq:primitive}: normalised i.i.d.\ rows stay a fixed distance from the
uniform row, since $\sum_kU_{ak}\sim m/2$ makes
$\delta_m\to\E|2U-1|=\tfrac12$. Incoming influence and stationary weights
nonetheless become increasingly balanced, and the numerical mixing deviation
$\max_a|\sum_jA_{aj}c_j-\bar c|$ falls with $m$ while $\bth'(0)$ approaches the
mean-field value $\mathrm{CV}^{2}(\vec t)$; what does this is not convergence of
rows but a law of large numbers, the weights being drawn independently of the
types. We make no formal convergence claim for the ensemble. Nothing in this
licenses a claim about networks in general --- least of all because the
independence it relies on is exactly what homophily destroys. Real influence structures carry
communities, hubs and correlation between characteristics and position, and
Remark~\ref{rem:blocks} shows that reversal survives at any population size once
such structure is present.
A supplement reports the frequency of both the cumulative event
$\bth(0.5)<\bth(0)$ and the marginal event $\bth'(0.5)<0$ over $20{,}000$ draws
at each of ten combinations of population size and dispersion, together with the
generators, the seeds and the analytic derivative used. Reversal is already
below one per cent at $m=6$ and indistinguishable from zero at $m\geq10$ in that
ensemble.

\begin{remark}[Reversal at any population size]
\label{rem:blocks}
The decay reported in the supplement reflects mixing, not size. Because
\eqref{eq:system} is local, a block-diagonal $A$ decouples: each block solves
its own system and $\bth=\big(m^{-1}\sum_ac_a\big)^{-1}$ is the reciprocal of
the population-average capacity. Replicating an identical block leaves that
average, and hence $\bth$, unchanged. Taking $k$ disjoint copies of the network
of Proposition~\ref{prop:reversal} therefore gives $\bth(\lambda)$ independent
of $k$, so that $\bth(1)=0.837209<6/7=\bth(0)$ at $m=3k$ for every $k$; we verify
this numerically up to $m=3000$. The construction is robust to weak ties across
blocks: mixing each row with a fraction $\epsilon$ of the uniform complete
network preserves the reversal at $m=30$ for $\epsilon\leq0.1$
($\bth(1)=0.8526$ at $\epsilon=0.1$) and destroys it by $\epsilon=0.2$
($\bth(1)=0.8681$). The simulation should therefore be read as describing a
dense ensemble in which local benchmarks approach a common average, and not as
a claim about large networks in general: block structure, or any persistent
correlation between type and network position, sustains the reversal at
arbitrary population size.

Note the contrast with Proposition~\ref{prop:homophily}, which uses the same
block structure to the opposite end. There net risk aversion is constant
\emph{within} blocks and the mechanism vanishes; here it is heterogeneous
within blocks and absent \emph{across} them, and the mechanism survives at any
scale. What separates the two is where the heterogeneity sits relative to the
partition, not the topology.
\end{remark}

\subsection{Perfect homophily blocks the mechanism}

\begin{proposition}
\label{prop:homophily}
Partition the investors into groups such that $A_{aj}=0$ whenever $a$ and $j$
belong to different groups, and suppose $\theta_a$ is constant within each
group. Then $c_a=1/\theta_a$ solves \eqref{eq:system} for every profile
$(\lambda_1,\dots,\lambda_m)$, and $\bth$ equals the harmonic mean of $\theta$
regardless of the intensity of comparison.
\end{proposition}

\begin{proof}
Row $a$ of \eqref{eq:system} reads
$(\theta_a+\lambda_a)c_a-\lambda_a\sum_jA_{aj}c_j=1$. If $c_j=1/\theta_j$ and
$\theta_j=\theta_a$ for every $j$ with $A_{aj}>0$, the left-hand side is
$(\theta_a+\lambda_a)/\theta_a-\lambda_a/\theta_a=1$. Uniqueness gives the
claim, and $\bth=m/\sum_a(1/\theta_a)$ is the harmonic mean of $\theta$.
\end{proof}

Theorem~\ref{thm:intense} explains why. Each group is a recurrent class, so
$r>1$, and $\theta$ is constant on $C_k$, whence
$\vec\pi^{(k)\top}\vec\theta=\theta_k$ whatever the stationary distribution
inside the class. With every investor in some class, $\mathbf 1^{\top}\vec h_k$
is the size $m_k$ of the class and \eqref{eq:intense} reads
$m/\sum_km_k/\theta_k=H(\theta)$. The two limits therefore coincide, which is
the strongest possible form of neutrality of the aggregate: not that $\bth$
returns to $H(\theta)$, but that it never leaves. With two equally sized groups
at $\theta\in\{0.5,2\}$, $\bth=0.8$ at $\lambda=0$, at $\lambda=1$ and at
$\lambda=1000$; the arithmetic mean $1.25$ that Proposition~\ref{prop:mf} would
predict is never approached.

The hypothesis is exact segregation together with exact constancy of $\theta$
within blocks, and both are needed: if $\theta$ varies inside a class then
$\vec\pi^{(k)\top}\vec\theta\neq\theta_k$ and \eqref{eq:intense} moves, as
Remark~\ref{rem:blocks} exploits. So the correct statement is that
\emph{perfect homophily by net risk-aversion type} suppresses the mechanism, not
that homophily does. It is worth knowing because assortativity on wealth and on
financial sophistication is a documented feature of social networks
\cite{mcpherson,bailey}.

%=====================================================================

%=====================================================================
\section{Conclusion}

Relative wealth concerns in a CARA economy aggregate through the influence
network in a way that two limits organise. The marginal effect at
zero comparison is a quadratic form in risk tolerances, which is a Dirichlet
form --- and so non-negative --- exactly when attention is balanced. The limit
under intense comparison is a stationary average of individual coefficients,
weighted by influence rather than by population. They anchor the two ends of the family rather than tracing the path between
them, but between them they account for the mean-field interpolation from the
harmonic to the arithmetic mean, for non-monotonicity and sign reversal, and
for the complete inertness of the aggregate under perfect homophily by type. What decides which
case obtains is the correlation between net risk aversion and stationary
influence.

Because the economy is closed, these are statements about prices and not only
about a conditional relation. Equilibrium prices are
$\vec p=(\vec\mu-\theta_MV\vec\nu)/R_f$ and the market Sharpe ratio is
$\bth\sqrt{\vec\nu^{\top}V\vec\nu}/m$, both explicit, because CARA demands carry
no wealth effects and $\theta_M$ therefore does not depend on prices. A network that lowers $\bth$ compresses the market risk premium and the market
Sharpe ratio. Its effect on an individual price has the sign of
$(V\vec\nu)_i$, so assets that covary positively with the market payoff rise
while hedges fall; and under the zero-net-supply closure, where
$\vec\nu^{\top}\vec p$ is pinned to aggregate initial wealth, the riskless rate
rises and value is redistributed across assets rather than created.

One boundary should be kept in view. Here $\theta_a$ is risk aversion net of the
comparison, not the coefficient of absolute risk aversion in own wealth, so
``common risk aversion'' in Theorem~\ref{thm:neutral} means a common value of
the product $\alpha_a(1-\gamma_a)$, along which stronger comparison motives come
paired with greater own-wealth risk aversion.

The practical reading is that a representative-agent treatment of relative
concerns reproduces the mean-field case and is safe only where influence is
close to evenly spread. Where attention is unbalanced --- where some investors
are watched much more than they watch --- it can err in sign. Two questions
follow naturally. The first is quantitative: how far can
$\vec\pi^{\top}\vec\theta$ fall below $H(\theta)$ for networks with realistic
degree distributions and realistic correlation between risk aversion and
centrality? The second is empirical: stationary influence is in principle
measurable, and Theorem~\ref{thm:intense} converts it into a prediction about
the aggregate.

%=====================================================================
\section*{Replication}
\label{sec:replication}

All numerical results are reproducible from the repository
%
\begin{center}
\url{https://github.com/renatovicente/price-of-risk-aggregation}
\end{center}
%
The module \texttt{scripts/network.py} solves \eqref{eq:system};
\texttt{scripts/neutrality.py} verifies Theorem~\ref{thm:neutral} and
Proposition~\ref{prop:mf}; \texttt{scripts/sign\_reversal.py} produces
the table in the supplement, including the network and type generators and
the seeds; \texttt{scripts/block\_reversal.py} produces the block replication and
the mixing robustness of Remark~\ref{rem:blocks}; and
\texttt{scripts/characterization.py} checks Theorems~\ref{thm:deriv} and
\ref{thm:intense}, the Dirichlet identity \eqref{eq:dirichlet}, the bound of
Lemma~\ref{lem:bound} and each of the three cases of \eqref{eq:intense}; and
\texttt{scripts/closed\_economy.py} checks the first-order conditions
\eqref{eq:foc} against direct numerical optimisation, market clearing, the price
formula \eqref{eq:price}, the recovery of \eqref{eq:clear} and the Sharpe
identity \eqref{eq:sharpe}. The scripts depend only on
\texttt{numpy} and \texttt{scipy}, read no data and draw from fixed seeds;
\texttt{./reproduce.sh} runs all six in about fifteen seconds.

%=====================================================================
\begin{thebibliography}{99}

\bibitem{abel} A. B. Abel,
\emph{Asset prices under habit formation and catching up with the Joneses},
American Economic Review Papers and Proceedings 80 (1990) 38--42.

\bibitem{ballester} C. Ballester, A. Calv\'o-Armengol, Y. Zenou,
\emph{Who's who in networks. Wanted: the key player},
Econometrica 74 (2006) 1403--1417.

\bibitem{campbell} J. Y. Campbell, J. H. Cochrane,
\emph{By force of habit: a consumption-based explanation of aggregate stock
market behavior}, Journal of Political Economy 107 (1999) 205--251.

\bibitem{chankogan} Y. L. Chan, L. Kogan,
\emph{Catching up with the Joneses: heterogeneous preferences and the dynamics
of asset prices}, Journal of Political Economy 110 (2002) 1255--1285.

\bibitem{degennaro} L. De Gennaro Aquino, E. G. De Giorgi, Y. Lou, M. S. Strub,
\emph{Investment and asset pricing with relative wealth concerns and multiple
risky assets},
The North American Journal of Economics and Finance 84 (2026) 102628.

\bibitem{bailey} M. Bailey, R. Cao, T. Kuchler, J. Stroebel,
\emph{The economic effects of social networks: evidence from the housing
market}, Journal of Political Economy 126 (2018) 2224--2276.

\bibitem{bermanplemmons} A. Berman, R. J. Plemmons,
\emph{Nonnegative Matrices in the Mathematical Sciences},
SIAM, Philadelphia, 1994.

\bibitem{demarzo04} P. DeMarzo, R. Kaniel, I. Kremer,
\emph{Diversification as a public good: community effects in portfolio choice},
Journal of Finance 59 (2004) 1677--1716.

\bibitem{dupor} B. Dupor, W.-F. Liu,
\emph{Jealousy and equilibrium overconsumption},
American Economic Review 93 (2003) 423--428.

\bibitem{gali} J. Gal\'i,
\emph{Keeping up with the Joneses: consumption externalities, portfolio choice,
and asset prices},
Journal of Money, Credit and Banking 26 (1994) 1--8.

\bibitem{gollier} C. Gollier,
\emph{Wealth inequality and asset pricing},
Review of Economic Studies 68 (2001) 181--203.

\bibitem{gomez} J.-P. G\'omez,
\emph{The impact of keeping up with the Joneses behavior on asset prices and
portfolio choice},
Finance Research Letters 4 (2007) 95--103.

\bibitem{kraft} H. Kraft, A. Meyer-Wehmann, F. T. Seifried,
\emph{Dynamic asset allocation with relative wealth concerns in incomplete
markets},
Journal of Economic Dynamics and Control 113 (2020) 103857.

\bibitem{lacker} D. Lacker, T. Zariphopoulou,
\emph{Mean field and $n$-agent games for optimal investment under relative
performance criteria},
Mathematical Finance 29 (2019) 1003--1038.

\bibitem{mcpherson} M. McPherson, L. Smith-Lovin, J. M. Cook,
\emph{Birds of a feather: homophily in social networks},
Annual Review of Sociology 27 (2001) 415--444.

\bibitem{panageas} S. Panageas,
\emph{The implications of heterogeneity and inequality for asset pricing},
NBER Working Paper 26974 (2020).

\end{thebibliography}
